% Generated draft: requires mathematical and source review.
\documentclass[11pt]{article}
\usepackage[margin=25mm]{geometry}
\usepackage{fontspec}
\setmainfont{Latin Modern Roman}
\usepackage{amsmath,amssymb,mathtools,unicode-math}
\setmathfont{Latin Modern Math}
\usepackage{enumitem,needspace,xurl,hyperref}
\setlength{\parindent}{0pt}
\setlength{\parskip}{5pt}
\setlength{\emergencystretch}{3em}
\newcommand{\N}{\mathbb{N}}
\newcommand{\Z}{\mathbb{Z}}
\newcommand{\Q}{\mathbb{Q}}
\newcommand{\R}{\mathbb{R}}
\newcommand{\C}{\mathbb{C}}
\newcommand{\F}{\mathbb{F}}
\newcommand{\A}{\mathbb{A}}
\newcommand{\PP}{\mathbb{P}}
\newcommand{\E}{\mathbb{E}}
\newcommand{\eps}{\varepsilon}
\newcommand{\dd}{\,\mathrm{d}}
\begin{document}
\section*{UnsolvedMath ID: 30004076}
\textbf{Source problem code:} catalogue entry
\subsection{Definitions and mathematical statement}
Let \(K\) be a field and let \(S=K[x_0,\ldots,x_n]\), with \(\deg x_i=1\), be a standard graded polynomial ring. Let \(I\subseteq S\) be a homogeneous ideal, put \(R=S/I\), and write \(X=\operatorname{Proj}(R)\). For a monomial order \(\prec\), the initial ideal \(\operatorname{in}_{\prec}(I)\) is generated by the leading monomials of elements of \(I\). It is square-free if it admits monomial generators in which every variable has exponent at most one.

Let \(\mathfrak m=R_+\) and \(d=\dim R\). The ring \(R\) is Cohen--Macaulay if \(\operatorname{depth}_{\mathfrak m}R=d\). Its graded \(a\)-invariant is
\[
 a(R)=\sup\{j\in\Z:H^d_{\mathfrak m}(R)_j\ne0\},
\]
where \(H^d_{\mathfrak m}(R)\) denotes graded local cohomology and the supremum of the empty set is \(-\infty\).

\textbf{Question.} Suppose that \(X\) is a nonsingular projective variety, as stipulated in the source, and that there exists a monomial order \(\prec\) for which \(\operatorname{in}_{\prec}(I)\) is square-free. Must
\[
 \operatorname{depth}_{\mathfrak m}R=\dim R
 \qquad\text{and}\qquad a(R)<0
\]
hold for every such \(K,S,I\)? No restriction on the characteristic of \(K\) or on the monomial order is stated in the supplied formulation.

\textbf{Curve formulation.} Under the same hypotheses, if \(\dim X=1\), must \(X\) have genus zero? For a smooth projective geometrically integral curve over \(K\), genus means
\[
 g(X)=\dim_K H^1(X,\mathcal O_X).
\]
The record does not specify geometric integrality or the intended genus convention outside that setting. Its original statement calls the curve question equivalent to the ring-theoretic conclusion; this claimed equivalence is retained as a source qualification requiring review.

\textbf{Unresolved hypotheses.} The source does not specify whether nonsingularity means regularity over \(K\) or smoothness over \(K\), whether ``variety'' requires geometric integrality, or whether \(I\) is intended to be prime or saturated. The expression \(X=\operatorname{Proj}(S/I)\) records the scheme-theoretic reading of ``defines''; these additional conventions have not been imposed.
\subsection{Short English statement}
The question asks whether a homogeneous coordinate ring defining a nonsingular projective variety and admitting a square-free initial ideal must be Cohen--Macaulay with negative \(a\)-invariant. It also asks whether every curve satisfying these hypotheses has genus zero. The supplied dated assessment reports results for degree-reverse-lexicographic order and for curves, while retaining an open label for the general assertion. Field, geometric, and equivalence conventions require review.
\subsection{Sources}
\begin{itemize}
\item Local UnsolvedMath registry, permanent ID \(30004076\), source code \texttt{catalogue entry}, title \emph{Smooth Varieties with Squarefree Gr\"obner Degenerations}, in the Oberwolfach Reports Open Problems collection. The supplied proposed year is 2019; no individual proposer is named.
\item Supplied report citation: \emph{Singularities and Homological Aspects of Commutative Algebra} (2020), Oberwolfach Reports, pp.~407--464, DOI \texttt{10.4171/owr/2019/6}. \url{https://doi.org/10.4171/owr/2019/6}. The background describes Oberwolfach Reports as documenting mathematical workshops and their open-problem discussions. The literature-source list labels this reference ``Oberwolfach Report 6/2019.''
\item Supplied literature reference: Huang--Tarasova--Varbaro--Witt, \url{https://arxiv.org/abs/2512.00584}. No title or theorem identifier is supplied.
\item Literature assessment checked at \texttt{2026-08-22}: the supplied curation labels the problem open and states that the conjecture is proved for degree-reverse-lexicographic order, that the curve case has since been settled, and that the full assertion for arbitrary square-free initial ideals remains open. This is supplied dated curation derived from the OWR material, not an independently verified literature assessment. The background explicitly requests independent verification of the time-sensitive status and cited claims.
\end{itemize}
\end{document}
